\clearpage
\fancyhead[L]{\small DURCHMESSER DER SONNE.}
\hypertarget{native-524-lower}{}\label{native-524-lower}
\begin{center}\small Göttingische gelehrte Anzeigen. Stück 113. Seite 1121..1126. 1809 Juli 17.\\[3pt]\rule{.72\linewidth}{.3pt}\end{center}

In einem Schreiben an unsern Hrn. Prof. \textsc{Gauss} vom 8. Juni hat Hr. \textsc{von Lindenau} die \emph{Resultate von Untersuchungen über den Durchmesser der Sonne} mitgetheilt, welche eine besondere Aufmerksamkeit verdienen. Hr. \textsc{von Lindenau} wurde zu diesen Untersuchungen veranlasst, als er bei dem Versuche, mit Hülfe einer ziemlich beträchtlichen Anzahl auf der Seeberger Sternwarte beobachteter Culminationen der Sonne den Sonnendurchmesser aus der Dauer der Durchgänge zu bestimmen, Unterschiede fand, welche er blos Beobachtungsfehlern zuzuschreiben Bedenken trug, da sie sehr deutlich einer bestimmten \emph{Periode} zu folgen schienen. Er entschloss sich also, zu diesem Behuf die sämmtlichen Greenwicher Beobachtungen von 1750 bis 1786 zu discutiren, und die Resultate dieser Untersuchung sind in der That sehr merkwürdig. Hr. \textsc{von Lindenau} fasste zuvörderst die Resultate für die durchgehends auf die mittlere Distanz der Erde von der Sonne reducirten Durchmesser der Sonne nach den verschiedenen einzelnen Jahren zusammen, wie folgendes Tableau zeigt:

\SourceLeaf{525}\begin{center}\small
\setlength{\tabcolsep}{3pt}
\begin{tabular}{c|c|c||c|c|c}
\shortstack{Jahr der\\Beobach-\\tung}&\shortstack{Mittleres Resul-\\tat für den Son-\\nenhalbmesser in\\der mittlern Ent-\\fernung}&\shortstack{Anzahl der\\Beobach-\\tungen}&\shortstack{Jahr der\\Beobach-\\tung}&\shortstack{Mittleres Resul-\\tat für den Son-\\nenhalbmesser in\\der mittlern Ent-\\fernung}&\shortstack{Anzahl der\\Beobach-\\tungen}\\ \hline
1750 & $962''63$ & 37 & 1773 & $961''42$ & 56 \\
1751 & $961.82$ & 42 & 1774 & $961.22$ & 32 \\
1752 & $961.61$ & 61 & 1775 & $960.86$ & 62 \\
1753 & $961.81$ & 53 & 1776 & $960.75$ & 53 \\
1754 & $962.30$ & 71 & 1777 & $960.45$ & 81 \\
1755 & $961.81$ & 46 & 1778 & $960.36$ & 71 \\
1765 & $962.44$ & 65 & 1779 & $960.92$ & 41 \\
1766 & $962.70$ & 46 & 1780 & $960.80$ & 37 \\
1767 & $961.57$ & 77 & 1781 & $960.09$ & 50 \\
1768 & $961.53$ & 72 & 1782 & $960.26$ & 37 \\
1769 & $961.23$ & 65 & 1783 & $959.84$ & 42 \\
1770 & $961.66$ & 47 & 1785 & $959.93$ & 65 \\
1771 & $961.82$ & 64 & 1786 & $959.65$ & 55 \\
1772 & $961.84$ & 56 & & & \\
\end{tabular}\end{center}

Es ist sehr auffallend, dass hier der Sonnenhalbmesser stufenweise immer abnimmt: inzwischen ist doch eine in so kurzer Zeit so beträchtliche Verminderung des Sonnenkörpers selbst, die für 3 Secunden Verminderung des Sonnenhalbmessers beinahe ein Hunderttheil des ganzen körperlichen Inhalts betragen würde, durchaus nicht denkbar, und man kann daher wohl den Grund dieses Phänomen nirgends anders suchen, als im Fernrohre, oder in dem Auge des Beobachters, oder in der Beobachtungsmanier. — Eben so auffallende Resultate gab die Zusammenstellung der einzelnen Bestimmungen nach den verschiedenen Jahreszeiten. Es zeigten sich hier, diese ganze Reihe von Jahren hindurch, und nur mit wenigen Ausnahmen, unverkennbare, periodisch wiederkehrende, Aenderungen: die mittlern Werthe für die einzelnen Monate, durchgehends auf den 15. reducirt, waren folgende:

\begin{center}\begin{tabular}{l|c||l|c}
&\shortstack{Halbmesser\\der Sonne}&&\shortstack{Halbmesser\\der Sonne}\\ \hline
Januar & $960''17$ & Juli & $960''14$ \\
Februar & $961.16$ & August & $961.00$ \\
März & $961.52$ & September & $961.70$ \\
April & $961.22$ & October & $961.80$ \\
Mai & $961.20$ & November & $961.16$ \\
Juni & $960.00$ & December & $960.43$ \\
\end{tabular}\end{center}

Es zeigten sich hier zwei Maxima, die man auf den Anfang Aprils und Octobers, und zwei Minima, die man auf den Anfang Januars und Julis setzen kann: der Unterschied der Extreme kann für den Halbmesser zu anderthalb, also für den Durchmesser zu drei Secunden angeschlagen werden. Es dürfte Manchem freilich etwas misslich scheinen, über so kleine Grössen durch \emph{Zeitbeobachtungen}, die einzeln genommen, oft viel grössere Differenzen zeigen, entscheiden zu wollen; von der andern Seite aber spricht die grosse Genauigkeit, mit welcher in Greenwich von \textsc{Bradley} und \textsc{Maskelyne} immer beobachtet worden ist, die so ungemein grosse Anzahl der Beobachtungen, und die Regularität in den Aenderungen der Resultate wieder zu sehr dagegen, jene Unterschiede blos für zufällig und für Folgen von Beobachtungsfehlern zu halten. Man könnte vielleicht zuerst darauf fallen, sie dem Einfluss der Jahreszeiten auf die Atmosphäre, also einer Modification der Irradiation, oder einem Einflusse auf die Empfindlichkeit des Auges, zuzuschreiben: allein dann müssten doch allem Anschein nach auf die Zeiten der grössten Wärme und Kälte die beiden verschiedenen Extreme fallen, anstatt dass die Beobach\SourceLeaf{518}tungen für beide Zeiten einerlei Durchmesser geben. Sehr ungezwungen würde sich hingegen das Phänomen durch eine \emph{elliptische Gestalt des Sonnenkörpers} erklären lassen. Nimmt man an, dass die Sonne ein durch Umdrehung um ihre Rotationsaxe erzeugtes Ellipsoid sei, und bezeichnet den Halbmesser des Sonnenäquators mit $a$, die halbe Axe mit $b$; bezeichnet man ferner die Neigung des Sonnenäquators gegen den Erdäquator mit $i$, die Rectascension des aufsteigenden Knotens jener Ebene auf dieser mit $N$, und die Rectascension der Sonne mit $A$, endlich den scheinbaren Halbmesser einer Kugel vom Halbmesser $a$, aus der Entfernung $1$ gesehen, mit $r$, so zeigt eine einfache Entwickelung, dass der aus der beobachteten Culminationszeit abgeleitete und auf die mittlere Entfernung reducirte Halbmesser sein werde

\[ r\left(1-\tfrac{1}{2}\left(1-\frac{bb}{aa}\right)\cos(A-N)^2\sin i^2\right) \]

Aus der von frühern Astronomen bestimmten Lage des Sonnenäquators, dessen aufsteigenden Knoten auf der Ekliptik in $78^\circ$ Länge, und die Neigung gegen die Ekliptik zu $7\tfrac{1}{2}$ Grad angenommen wird, folgt $i = 26^\circ \, 2'$, $N = 16^\circ \, 55'$. Die äussersten Werthe des Halbmessers werden also Statt haben, wenn $A-N$ entweder $0$, oder $90^\circ$, oder $180^\circ$, oder $270^\circ$ ist; im ersten und dritten Fall, d. i. den 8. April und 12. October, wird jener Halbmesser sein $= r\left(1-\left(1-\dfrac{bb}{aa}\right).\,0{,}09631\right)$, im zweiten und vierten Fall hingegen, also am 8. Juli und 6. Januar, wird man den Halbmesser $= r$ finden: jener wird der kleinere oder grössere sein, je nachdem die Sonne ein abgeplattetes oder längliches Sphäroid ist. Sehr sonderbar ist es, dass die beobachteten Phänomene bei dieser Erklärungsart die letztere Hypothese erfordern. In der That werden dieselben sehr gut dargestellt, wenn man $\dfrac{bb}{aa}-1 = \tfrac{1}{58}$, oder $\dfrac{b}{a} = \tfrac{17}{116}$ setzt oder die Aequatorialabplattung des Sonnenkörpers $\tfrac{1}{117}$. Hr. \textsc{von Lindenau} findet etwas weniger, nemlich $\tfrac{1}{140}$. — Um zu sehen, was directe Messungen über diese Gestalt entscheiden würden, hat Hr. \textsc{von Lindenau} nun auch die verticalen Sonnendurchmesser aus den Unterschieden der Zenithdistanzen der beiden Ränder mit Hülfe einer eben so grossen Anzahl von Beobachtungen abgeleitet. Folgendes sind hiervon die Resultate:

\begin{center}\begin{tabular}{c|c||c|c}
Jahr & \shortstack{Mittlerer Son-\\nenhalbmesser} & Jahr & \shortstack{Mittlerer Son-\\nenhalbmesser} \\ \hline
1765 & $963''81$ & 1775 & $963''15$ \\
1766 & $964.28$ & 1776 & $962.71$ \\
1767 & $963.88$ & 1777 & $961.74$ \\
1768 & $963.40$ & 1778 & $962.47$ \\
1769 & $964.07$ & 1779 & $962.85$ \\
1770 & $963.31$ & 1780 & $962.63$ \\
1771 & $962.89$ & 1781 & $961.94$ \\
1772 & $963.36$ & 1782 & $961.84$ \\
1773 & $963.16$ & 1783 & $961.52$ \\
1774 & $963.46$ & & \\
\end{tabular}\end{center}

Auch hier zeigt sich also eben so, wie bei den horizontalen Durchmessern, eine successive Verminderung, und durchgehends sind die verticalen um einige Secunden (im Mittel um $3''65$) grösser, als jene. Hr. \textsc{von Lindenau} berechnete hieraus eine Aequatorialabplattung von $\tfrac{1}{288}$; Hr. Prof. \textsc{Gauss} findet noch etwas weniger, nemlich $\tfrac{1}{420}$. Man sieht also, dass sowol die Vergleichung der zu verschiedenen Jahrszeiten gemessenen horizontalen Sonnendurchmesser unter sich, als die Vergleichung der mittlern horizontalen mit den mittlern verticalen, in so fern man sich auf die Schärfe der Messungen selbst verlassen kann, sich vereinigen, der Sonne eine längliche Gestalt zu geben, obwohl die Verhältnisse der Ellipticität nicht harmoniren. Je weniger, bei einer solchen Gestalt des Sonnenkörpers, die Möglichkeit eines Zustandes von Gleichgewicht auf seiner Oberfläche begriffen wird, um so begieriger wird \SourceLeaf{519}man sein, das nähere Detail dieser interessanten Untersuchung des Hrn. \textsc{von Lindenau} bald kennen zu lernen.

Wir benutzen noch diese Gelegenheit, ein von dem Herrn \textsc{von Wisniewski}, ausserordentlichem Mitgliede der Petersburgischen Akademie der Wissenschaften, aus Astracan an den Hrn. Dr. \textsc{Olbers} eingesandtes und von diesem uns mitgetheiltes Verzeichniss einiger noch in diesem Jahre vorfallenden und in \textsc{Bode}'s Astronomischem Jahrbuche nicht angezeigten Sternbedeckungen vom Monde hier mitzutheilen, um die Beobachter auf dieselben aufmerksam zu machen, und so vielleicht zu den von Hrn. \textsc{von Wisniewski} im Russischen Reiche anzustellenden Beobachtungen einige correspondirende zu verschaffen. Das Verzeichniss ist für den Horizont eines Orts unter der Breite $52^{\circ}\,31'$, und der östlichen Länge von Greenwich $2^{\mathrm u}\,0^{\mathrm m}$ berechnet. — —

Mehrere dieser Bedeckungen sind bereits in der Connoissance des tems angezeigt.
